(1) 球对称 Helmholtz 方程为 $(r^2p_r)_r+k^2r^2p=0$。令 $u=rp$，得到 $u''+k^2u=0$，向外辐射解为 $\widetilde p(r)=A e^{ikr}/r$，真实压力取 $\Re\{\widetilde p e^{-i\omega t}\}$。只在远离点源局部反应区时将其视为近似传播球面波。
(2) 无耗散介质中声压幅值 $p_a=|A|/r$。远区速度幅值 $v_a\simeq p_a/(\rho c)$，平均声强 $I=|p_a|^2/(2\rho c)=|A|^2/(2\rho c r^2)$。穿过半径 $r$ 球面的总功率 $\mathcal P=4\pi r^2I=2\pi|A|^2/(\rho c)$ 与 $r$ 无关，印证球面扩散是几何原因而非耗散。
(3) 有吸收时复波数为 $k=k_r+i\alpha$，故 $e^{ikr}=e^{ik_rr-\alpha r}$，压力幅度 $p_a(r)=|A|e^{-\alpha r}/r$。在弱吸收且局部关系近似保持的条件下，声强 $I(r)\simeq|A|^2e^{-2\alpha r}/(2\rho cr^2)$；强度衰减指数有因模平方而产生的因子二。
(4) 两个距离的声强比为 $I(r_2)/I(r_1)=(r_1/r_2)^2e^{-2\alpha(r_2-r_1)}$。前因子仅来自能量流过更大球面，即使 $\alpha=0$ 仍有；指数才是单位传播长度上的介质耗散。若实验将纯几何衰减误判成材料吸声，会高估衰减系数。
